\documentclass[12pt,a4paper]{article}
\usepackage[UTF8]{ctex}
\usepackage{amsmath,amsfonts,amssymb}
\usepackage{geometry}
\usepackage{fancyhdr}
\usepackage{xcolor}
\usepackage{tcolorbox}

% 保持与上一份文件一致的格式设置
\geometry{left=2cm,right=2cm,top=2.5cm,bottom=2.5cm}
\pagestyle{fancy}
\fancyhf{}
\fancyhead[L]{\textbf{高阶导数公式推导详解}}
\fancyhead[R]{自己备课用}
\fancyfoot[C]{\thepage}
\renewcommand{\headrulewidth}{0.4pt}
\setlength{\parindent}{0pt}
\setlength{\parskip}{0.8em}

% 定义颜色
\definecolor{derivecolor}{RGB}{0, 102, 102} % 深青色用于推导部分
\newtcolorbox{derivebox}{colback=derivecolor!5!white,colframe=derivecolor,title=\textbf{推导演示板书}}

\begin{document}

\section*{附录：常见 $n$ 阶导数公式的推导细节}

\textbf{授课建议}：讲解时，建议采用“计算前三阶 $\to$ 观察规律 $\to$ 归纳通式”的步骤。

\vspace{1em}

% --- 公式 1: 指数函数 ---
\subsection*{1. 指数函数 $(e^{ax})^{(n)}$ 的推导}

\begin{derivebox}
设 $y = e^{ax}$。
\begin{itemize}
    \item \textbf{一阶导}：$y' = (e^{ax})' = a \cdot e^{ax}$
    \item \textbf{二阶导}：$y'' = (a e^{ax})' = a \cdot (a e^{ax}) = a^2 e^{ax}$
    \item \textbf{三阶导}：$y''' = (a^2 e^{ax})' = a^2 \cdot (a e^{ax}) = a^3 e^{ax}$
\end{itemize}
\textbf{规律}：求导 $n$ 次，系数 $a$ 就乘 $n$ 次。
\[
\boxed{(e^{ax})^{(n)} = a^n e^{ax}}
\]
\end{derivebox}

\newpage
% --- 公式 2: 正弦函数 ---
\subsection*{2. 正弦函数 $(\sin ax)^{(n)}$ 的推导}
\textbf{难点}：如何将 $\cos$ 变回 $\sin$ 以寻找规律？利用诱导公式 $\cos \theta = \sin(\theta + \frac{\pi}{2})$。

\begin{derivebox}
设 $y = \sin ax$。
\begin{itemize}
    \item \textbf{一阶导}：
    \[ y' = a \cos ax = a \sin\left(ax + \dfrac{\pi}{2}\right) \]
    
    \item \textbf{二阶导}：
    \[ y'' = \left[ a \sin\left(ax + \dfrac{\pi}{2}\right) \right]' = a^2 \cos\left(ax + \dfrac{\pi}{2}\right) = a^2 \sin\left(ax + \dfrac{\pi}{2} + \dfrac{\pi}{2}\right) = a^2 \sin\left(ax + 2 \cdot \dfrac{\pi}{2}\right) \]
    
    \item \textbf{三阶导}：
    \[ y''' = a^3 \sin\left(ax + 3 \cdot \dfrac{\pi}{2}\right) \]
\end{itemize}
\textbf{规律}：求导一次，系数乘 $a$，相位增加 $\dfrac{\pi}{2}$。
\[
\boxed{(\sin ax)^{(n)} = a^n \sin\left(ax + \dfrac{n\pi}{2}\right)}
\]
\end{derivebox}

\newpage
% --- 公式 3: 余弦函数 ---
\subsection*{3. 余弦函数 $(\cos ax)^{(n)}$ 的推导}
\textbf{思路}：利用诱导公式 $\cos' x = -\sin x = \cos(x + \frac{\pi}{2})$。

\begin{derivebox}
设 $y = \cos ax$。
\begin{itemize}
    \item \textbf{一阶导}：
    \[ y' = -a \sin ax = a \cos\left(ax + \dfrac{\pi}{2}\right) \]
    
    \item \textbf{二阶导}：
    \[ y'' = \left[ a \cos\left(ax + \dfrac{\pi}{2}\right) \right]' = -a^2 \sin\left(ax + \dfrac{\pi}{2}\right) = a^2 \cos\left(ax + 2 \cdot \dfrac{\pi}{2}\right) \]
    
    \item \textbf{三阶导}：
    \[ y''' = a^3 \cos\left(ax + 3 \cdot \dfrac{\pi}{2}\right) \]
\end{itemize}
\textbf{规律}：与正弦函数完全一致，只是函数名保持为 $\cos$。
\[
\boxed{(\cos ax)^{(n)} = a^n \cos\left(ax + \dfrac{n\pi}{2}\right)}
\]
\end{derivebox}


\newpage

% --- 公式 4: 幂函数/分式 ---
\subsection*{4. 分式函数 $\left(\dfrac{1}{x+a}\right)^{(n)}$ 的推导}
\textbf{技巧}：写成负指数形式 $(x+a)^{-1}$ 更容易看清系数变化。

\begin{derivebox}
设 $y = (x+a)^{-1}$。
\begin{itemize}
    \item \textbf{一阶导}：
    \[ y' = (-1)(x+a)^{-2} = -\dfrac{1!}{(x+a)^2} \]
    
    \item \textbf{二阶导}：
    \[ y'' = (-1)(-2)(x+a)^{-3} = (-1)^2 \cdot (1 \cdot 2) (x+a)^{-3} = \dfrac{(-1)^2 2!}{(x+a)^3} \]
    
    \item \textbf{三阶导}：
    \[ y''' = (-1)^2 (2) \cdot (-3)(x+a)^{-4} = (-1)^3 \cdot (1 \cdot 2 \cdot 3) (x+a)^{-4} = \dfrac{(-1)^3 3!}{(x+a)^4} \]
\end{itemize}
\textbf{观察规律}：
\begin{enumerate}
    \item 符号：奇负偶正 $\to (-1)^n$
    \item 系数：$n!$
    \item 指数：分母次数为 $n+1$
\end{enumerate}
\[
\boxed{\left(\dfrac{1}{x+a}\right)^{(n)} = \dfrac{(-1)^n n!}{(x+a)^{n+1}}}
\]
\end{derivebox}

\newpage
% --- 公式 5: 对数函数 ---
\subsection*{5. 对数函数 $(\ln(x+a))^{(n)}$ 的推导}
\textbf{关键}：对数函数求导一次后，就变成了分式函数，所以它比分式函数“慢一拍”。

\begin{derivebox}
设 $y = \ln(x+a)$。
\begin{itemize}
    \item \textbf{一阶导}：
    \[ y' = \dfrac{1}{x+a} = (x+a)^{-1} \]
    \textit{(此时还没出现阶乘，相当于 $0!$)}
    
    \item \textbf{二阶导}：
    \[ y'' = \left[ (x+a)^{-1} \right]' = (-1)(x+a)^{-2} = -\dfrac{1!}{(x+a)^2} \]
    
    \item \textbf{三阶导}：
    \[ y''' = (-1)(-2)(x+a)^{-3} = \dfrac{2!}{(x+a)^3} \]
    
    \item \textbf{四阶导}：
    \[ y^{(4)} = (2!) \cdot (-3)(x+a)^{-4} = -\dfrac{3!}{(x+a)^4} \]
\end{itemize}
\textbf{观察规律}（注意 $n$ 与阶乘的关系）：
\begin{enumerate}
    \item 阶数是 $n$ 时，分子是 $(n-1)!$。
    \item 符号项由 $n-1$ 决定（一阶为正，二阶为负），即 $(-1)^{n-1}$。
\end{enumerate}
\[
\boxed{(\ln(x+a))^{(n)} = \dfrac{(-1)^{n-1} (n-1)!}{(x+a)^n}} \quad (n \geq 1)
\]
\end{derivebox}

\end{document}